% Conclusion and Phase 2 outlook (issue #7).

\section{Conclusion}
\label{sec:conclusion}

Phase~1 established a leakage-controlled protocol and feature-based baselines for
RUL prediction on C-MAPSS FD001. With 155 causal engineered features, the Random
Forest predicts the RUL of the 100 official test engines with an RMSE of 27.55
cycles and an MAE of 17.67, the lowest MAE of the six models; all six lie within
five RMSE cycles of each other. The overall numbers hide the main finding. Near
failure, where maintenance is decided, the engineered features roughly halve the
error of both tree models, to 7.52 and 5.62 cycles. Early in life, under the
uncapped target, they overestimate engines that still look healthy, and a single
such engine accounts for 92\% of the Random Forest's NASA score. Capping the
training target at 125 cycles removes much of this error and reverses the
raw-versus-engineered comparison, so conclusions about feature engineering
depend on how RUL is defined.

These results have limits. Only FD001, with one operating condition and one fault
mode, was evaluated, so regime normalization has not yet been tested where it
matters. Each model was tuned over a small grid with three folds, the XGBoost
learning rate was not tuned, and cross-validation scores every training cycle
rather than last cycles like the test set. Differences between models come from
one test set of 100 engines without confidence intervals, and the capped-target
result covers only the Random Forest.

Phase~2 keeps the splits, metrics and last-cycle scoring fixed and replaces the
fixed rolling summaries with sequence models, an LSTM, a temporal CNN and a
Transformer encoder, trained on sensor windows built inside engine boundaries.
A hybrid model will fuse a sequence encoder with the engineered-feature branch,
targeting the gap identified in Section~\ref{sec:introduction}. Ablations will
cover the capped target with the cap chosen on training data, window length,
feature groups and regime normalization on FD002 and FD004, and validation will
add repeated seeds, bootstrap confidence intervals and paired per-engine tests
using the saved predictions. The bar to beat is the Phase~1 method's test RMSE of
27.55 cycles (26.15 with raw inputs) and its late-stage RMSE of 6 to 8 cycles.
